\documentclass[a4paper,twoside]{article}

\usepackage{morewrites}

\usepackage[svgnames,dvipsnames,x11names]{xcolor}
\usepackage{graphicx}
\usepackage[english]{babel}
\usepackage[colorlinks=true, linkcolor=NavyBlue, urlcolor=NavyBlue, citecolor=NavyBlue]{hyperref}
\usepackage[a4paper]{geometry}
\usepackage{ts-fonts}
\usepackage{float}
\usepackage{seqsplit}

\usepackage{lastpage}
\usepackage{refcount}
\makeatletter
\def\ps@plain{
  \let\@oddhead\@empty
  \let\@evenhead\@empty
  \def\@oddfoot{
    \hfil
    \ifnum\getpagerefnumber{LastPage}>1\relax
      \thepage
    \fi
    \hfil}
  \let\@evenfoot\@oddfoot
}
\makeatother

\title{Glossary}
\author{}\date{}

\begin{document}

\maketitle

\thispagestyle{plain}
In online documentation, a glossary doesn't make much sense because you can
search for terms directly and you have hyperlinks. However, in printed documents, a
glossary can be very useful to provide definitions of terms used in the text.

TeXSmith adds support for glossaries through the \texttt{glossary} extension, which
allows you to define glossary entries in your Markdown files and generate a
glossary section in the output document.

The fragment is automatically included when needed:

\begin{itemize}
\item{} If you use acronyms or abbreviations.
\item{} If you define glossary entries using the \texttt{glossary} directive.

\end{itemize}

\end{document}
